% Industry resume, Grover Aruquipa
% Compile with pdfLaTeX (also works on Overleaf). One column on purpose:
% applicant tracking systems read single-column text far more reliably than sidebars.
\documentclass[10pt,letterpaper]{article}
\usepackage[T1]{fontenc}
\usepackage[utf8]{inputenc}
\usepackage[scaled=0.95]{helvet}
\renewcommand{\familydefault}{\sfdefault}
\usepackage[top=1.1cm,bottom=1.0cm,left=1.5cm,right=1.5cm]{geometry}
\usepackage{enumitem}
\usepackage{titlesec}
\usepackage{xcolor}
\usepackage{microtype}
\usepackage[hidelinks]{hyperref}

\definecolor{ink}{HTML}{111111}
\definecolor{accent}{HTML}{0D2B61}
\definecolor{muted}{HTML}{555555}
\color{ink}
\pagestyle{empty}
\setlength{\parindent}{0pt}
\setlength{\parskip}{0pt}

\titleformat{\section}{\normalsize\bfseries\color{accent}}{}{0em}{\MakeUppercase}[\vspace{-6pt}{\color{accent}\rule{\linewidth}{0.6pt}}]
\titlespacing*{\section}{0pt}{7pt}{3pt}
\setlist[itemize]{leftmargin=1.1em,itemsep=0.5pt,topsep=1pt,parsep=0pt,label=\textbullet}

% \role{Title}{Organization}{Location}{Dates}
\newcommand{\role}[4]{\textbf{#1}, #2 \hfill {\color{muted}#3 \textbar\ #4}\par}

\begin{document}

%================ HEADER ================
{\LARGE\bfseries Grover Aruquipa}\par\vspace{2pt}
{\large Robotics Mechanical Engineer \textbar\ Mechanism Design and Multi-DoF Robotic Hardware}\par\vspace{3pt}
{\small Qu\'ebec, QC, Canada \ \textbullet\ \
\href{mailto:grover.aruquipa.9@gmail.com}{grover.aruquipa.9@gmail.com} \ \textbullet\ \
\href{https://www.linkedin.com/in/grover-aruquipa-996917126}{linkedin.com/in/grover-aruquipa-996917126} \ \textbullet\ \
\href{https://groveraruquipa.github.io}{groveraruquipa.github.io}}\par
% TODO: add a current phone number on the line above if you want it on the CV.

%================ SUMMARY ================
\section{Summary}
Mechanical engineer finishing a PhD (Universit\'e Laval, defense January 2027) on the design of non-traditional
multi-DoF robotic mechanisms. I have taken three robots with unusual kinematic structures from concept and CAD through
fabrication, wiring and real-time control to a working, validated prototype, including a compact coaxial wrist-gripper
that turns its tool without limit while holding it. Master's research in micro-manipulation at FEMTO-ST (CNRS).

%================ SKILLS ================
\section{Technical Skills}
\begin{itemize}
\item \textbf{Mechanism design:} parallel and hybrid mechanisms with unusual kinematic structures, coaxial and concentric
transmissions (hollow shafts, belts), revolute and prismatic joints, integrated grippers, vendor component selection.
\item \textbf{CAD, analysis and fabrication:} SolidWorks, Creo, ANSYS; machining, SLS and FDM printing, assembly and
commissioning of research prototypes.
\item \textbf{Electronics:} Altium (basic PCB design), motor and encoder wiring, component selection for multi-axis
drives.
\item \textbf{Kinematics and dynamics:} rigid-body kinematics and dynamics, screw theory and Lie groups, Jacobians,
singularity and workspace analysis, simulation in MATLAB and MuJoCo.
\item \textbf{Control and software:} real-time multi-axis control at 500\,Hz (OPAL-RT, Simulink Real-Time), C++ (real
time), Python, MATLAB, ROS~2, MoveIt, Pinocchio.
\end{itemize}

%================ EXPERIENCE ================
\section{Experience}
\role{Doctoral Researcher}{Robotics Laboratory, Universit\'e Laval}{Qu\'ebec, Canada}{2023 -- present}
\begin{itemize}
\item Designed, built and controlled three parallel manipulators (4, 7 and 8 actuators) with every motor
mounted on the base and the gripper driven through the mechanism, enabling unlimited rotation of the grasped object.
\item Implemented real-time multi-motor control at 500\,Hz; validated a 4-DoF coaxial wrist-gripper (0.25\,kg moving mass,
180\,mm footprint) over ten continuous tool revolutions with sub-degree tracking.
\item Derived closed-form inverse kinematics, Jacobians and singularity conditions, and developed a redundancy
strategy that keeps unlimited rotation free of collisions and singularities.
\item Led two graduate students from concept to working prototype on an 8-actuator (6+2)-DoF robot; joint paper
submitted to ICRA 2027.
\item First author of 3 journal papers (ASME J.\ Mechanisms and Robotics, Mechanism and Machine Theory) and 4 more
under review; reviewer for ASME J.\ Mechanisms and Robotics.
\end{itemize}

\vspace{3pt}
\role{Visiting Researcher}{Institute of Robotics, Johannes Kepler University Linz}{Linz, Austria}{Apr. -- May 2026}
\begin{itemize}
\item Developed grasp-quality indices for manipulators with integrated grasping (Mitacs Globalink Research Award).
\end{itemize}

\vspace{3pt}
\role{Research Intern}{FEMTO-ST (CNRS)}{Besan\c{c}on, France}{2022 -- 2023}
\begin{itemize}
\item Designed and modeled a redundant dexterous parallel robot with two unlimited rotations for micro-manipulation.
\item Wrote real-time control software in C++ for the robot experiments.
\end{itemize}

\vspace{3pt}
\textbf{Industry roles in vision and machine learning} \hfill {\color{muted}2020 -- 2021}\par
\begin{itemize}
\item \textbf{Computer Vision Engineer}, Sencon Srl (remote, 2021): computer vision for industrial anomaly detection.
\item \textbf{Graduate Research Intern}, Hitachi ABB Power Grids R\&D (Zurich, 2020): machine learning on sensor data from high-voltage circuit breakers.
\item \textbf{Data Scientist}, KTH Royal Institute of Technology (remote, 2021): data analysis for environmental impact assessment.
\end{itemize}

%================ EDUCATION ================
\section{Education}
\textbf{PhD, Mechanical Engineering}, Universit\'e Laval \hfill {\color{muted}Qu\'ebec, Canada \textbar\ 2023 -- Jan.\ 2027 (expected)}\par
{\small Kinematically redundant parallel robots with unlimited rotations and integrated grasping. Advisors: C.~Gosselin, P.~Lambert.}\par\vspace{2pt}
\textbf{MSc, Smart Integrated Systems}, Universit\'e Bourgogne--Franche-Comt\'e \hfill {\color{muted}Besan\c{c}on, France \textbar\ 2021 -- 2023}\par\vspace{2pt}
\textbf{BEng, Mechatronics Engineering}, Bolivian Catholic University \hfill {\color{muted}La Paz, Bolivia \textbar\ 2014 -- 2019}\par
{\small Thesis: reconfigurable Delta robot for product picking.}

%================ AWARDS ================
\section{Awards and Languages}
\begin{itemize}
\item 1st Place, Graduate Division, ASME Student Mechanisms and Robotics Design Competition (2025), team lead.
\item Finalist, MassRobotics competition, Boston, with one of my robot prototypes, team lead.
\item Mitacs Globalink Research Award (2026); 1st Place, Student Paper Competition, Argentine Robotics Conference (2019).
\item Languages: Spanish (native), English (full professional), French (intermediate).
\end{itemize}

\end{document}
